\documentclass[zihao=-4,a4paper]{ctexart}
\usepackage[margin=2.4cm]{geometry}
\usepackage{amsmath,amssymb,booktabs,graphicx,float,array,enumitem,xcolor}
\usepackage[most]{tcolorbox}
\usepackage[hidelinks]{hyperref}
\linespread{1.35}
\setlist{itemsep=2pt,topsep=4pt}
\ctexset{section={format=\Large\bfseries\color{blue!55!black}},subsection={format=\large\bfseries}}
\newtcolorbox{keybox}[1][本次要点]{colback=blue!3,colframe=blue!45!black,fonttitle=\bfseries,title=#1,boxrule=0.6pt,arc=1mm}
\newtcolorbox{notebox}{colback=gray!6,colframe=gray!60,boxrule=0.4pt,arc=1mm}
\newcommand{\fig}[3][0.95]{\begin{figure}[H]\centering\includegraphics[width=#1\linewidth]{#2}\caption{#3}\end{figure}}
\newcommand{\M}{M^{(K)}}

\title{\textbf{第 6 次组会汇报}\\[4pt]\Large 效率与外部方法对比}
\author{对应工作：126、127 号实验；论文当前版本}
\date{}

\begin{document}
\maketitle
\thispagestyle{empty}

\begin{keybox}
\begin{enumerate}[leftmargin=*]
  \item \textbf{统一计时}（同一张空闲 GPU、同一批集合）：最终估计器每个集合 48—64 ms，比串行重训快 241—906 倍；读出改为数学等价的快速实现，结果逐位相同，快约 2.8 倍。
  \item \textbf{与外部“免重训”方法对比}：均值代替、LazyVI、热启动、掩码代理模型直接输出、TabPFN v2。只有 TabPFN 的精度与本文相当，但\textbf{每集合慢 13—110 倍}；认证之后两者恢复的风险相同。
  \item 如实说明：TabPFN 的优势集中在训练行最少的 PJM，本文在那里有 0.03—0.04 的系统低估；它有时比重训的攻击器族还强，所以更适合作为攻击器族的一员。
\end{enumerate}
\end{keybox}

\section{统一计时（126 号）}

\subsection{先澄清“三种子”}

\begin{table}[H]\centering
\caption{各方法每个查询集合的计算量}
\small
\begin{tabular}{lccc}
\toprule
方法 & 每个查询集合的计算 & 读出次数 & 特征维度 \\
\midrule
旧 E（只预测目标 $\times3$） & 三个主干各自读出，预测取平均 & 3 & 300—326 \\
C2（重建 $\oplus$ 随机 $\times1$） & 两个主干特征拼接，读出一次 & 1 & 556—582 \\
\textbf{本文 C3（重建 $\oplus$ 随机 $\times3$）} & C2 做三遍，预测取平均 & 3 & 556—582 \\
三种子特征拼接（未采用） & 三个主干特征拼成一块，读出一次 & 1 & 812—838 \\
\bottomrule
\end{tabular}
\end{table}

论文里的“$\times3$”都是\textbf{三次独立读出再平均}，而不是把三个种子的特征拼在一起读出一次。后者反而更慢（41 ms 对 21 ms），认证也更差（0.916 对 0.953）。

\subsection{每个集合的耗时}

\fig{figures/图1.png}{126 号：各方法每集合耗时（灰：原实现；蓝：等价快速实现；红色虚线为串行重训约 30 秒 / 集合）\protect\newline{\footnotesize 原图：\texttt{DNN\_Aggresvation126/figures/fig1\_time\_per\_method.png}}}

\textbf{等价快速实现}：原实现在 5 折 $\times$ 6 个 $\lambda$ 中，每折都重新计算 Gram 矩阵，每个 $\lambda$ 单独求解。改为全体训练行的 $Z^\top Z$、$Z^\top y$ 只算一次，各折用“全体减去验证折”得到，6 个 $\lambda$ 合并为一次批量求解。只改计算顺序，\textbf{与原实现、与 125 号保存的估计值逐位相同}，读出快约 2.8 倍。

\fig{figures/图2.png}{126 号：耗时与精度的关系（左：全部方法；中：带主干的读出放大；右：认证前 1 比）\protect\newline{\footnotesize 原图：\texttt{DNN\_Aggresvation126/figures/fig2\_time\_vs\_accuracy.png}}}

\begin{table}[H]\centering
\caption{126 号：代表方法的耗时与精度（快速实现，五组数据）}
\small
\begin{tabular}{lccc}
\toprule
方法 & 每集合耗时 & $M^{(2)}$ 误差 & 认证前 1 \\
\midrule
共享输出头 $\times3$（不做逐集合读出） & 0.2 ms & 0.057 & 0.883 \\
随机主干读出 $\times1$ + 截断 & 5.7—8.0 ms & 0.021 & 0.958 \\
旧 E & 16.9—23.5 ms & 0.021 & 0.934 \\
\textbf{本文 C3} & \textbf{48.4—63.9 ms} & \textbf{0.015} & \textbf{0.961} \\
热启动微调 $\times1$（100 步） & 227—299 ms & $V$ 误差 10/10 个目标更差 & — \\
串行重训（三攻击器） & 14—56 s & 精确 & 精确 \\
\bottomrule
\end{tabular}
\end{table}

\begin{itemize}[leftmargin=*]
  \item 读出耗时基本由特征维度决定：约 300 维 6—8 ms，约 570 维 16—21 ms。
  \item 最省事的共享输出头快两个数量级，但 $M$ 误差约为本文的 4 倍。
  \item \textbf{随机主干 + 截断}是很便宜的备选：7 ms，$M$ 误差与旧 E 相当，但仍比本文差 40\%。
  \item 扫描一张完整的 $K=2$ 表，本文需要 1—7 分钟，串行重训需要 7—53 小时。一次性预训练每个重建主干约 35—46 秒。
\end{itemize}

\section{与外部“免重训”方法对比（127 号）}

\subsection{选了哪些方法}

论文中估计器的对比分成两块：\textbf{设计选择}（第 5 次讲的各种变体）和\textbf{与外部通用方法的对比}（本节）。外部方法都用默认设置、按原文实现，不做额外调优。

\begin{table}[H]\centering
\caption{127 号：外部 baseline}
\small
\begin{tabular}{lll}
\toprule
类别 & 方法 & 出处 \\
\midrule
全模型不重训 & 被移除字段以均值代替（Dropout） & Williamson 等；Sun \& Raskutti \\
全模型线性化 & LazyVI（参数处线性化 + 岭修正） & Gao 等，ICML 2022 \\
全模型热启动 & 热启动 + 早停 & Sun \& Raskutti，AISTATS 2025 \\
掩码代理模型 & 随机掩码网络直接输出 & Covert 等，JMLR 2021；FastSHAP \\
表格基础模型 & TabPFN v2（上下文推断，不训练） & Hollmann 等，Nature 2025 \\
\bottomrule
\end{tabular}
\end{table}

\subsection{结果}

\fig{figures/图3.png}{127 号：与外部方法的对比（左：$M^{(2)}$ 误差对耗时；中：认证前 1 比对耗时；右：同一批 200 个集合上的 $V$ 误差）\protect\newline{\footnotesize 原图：\texttt{DNN\_Aggresvation127/figures/fig1\_baselines.png}}}

\begin{table}[H]\centering
\caption{127 号：10 个目标平均（LazyVI 太慢，只在 200 个集合上算了 $V$ 误差 0.037）}
\footnotesize\setlength{\tabcolsep}{3pt}
\begin{tabular}{lccccc}
\toprule
方法 & $M^{(2)}$ 误差 & 排序 Spearman & 认证前 1 & 关键召回 $\tau=0.5$ / $0.7$ & 每集合耗时 \\
\midrule
串行重训（真值） & 0 & 1 & 1 & 1 / 1 & 14—56 s \\
逐集合线性回归 & 0.144 & 0.737 & 0.803 & 0.24 / 0.16 & 0.7—1.0 ms \\
均值代替（Dropout） & 0.243 & 0.534 & 0.295 & 0.03 / 0.00 & 0.1 ms \\
LazyVI（抽样 200 集合） & — & — & — & — & 1.5—9.1 s \\
热启动 + 早停 & 0.047 & 0.932 & 0.910 & 0.67 / 0.47 & 176—181 ms \\
掩码代理模型直接输出 & 0.057 & 0.930 & 0.883 & 0.55 / 0.47 & 0.2 ms \\
\textbf{TabPFN v2} & \textbf{0.013} & \textbf{0.979} & 0.959 & 0.88 / \textbf{1.00} & 0.8—7.0 s \\
\textbf{本文} & 0.015 & 0.968 & \textbf{0.961} & \textbf{0.91} / 0.79 & \textbf{48—64 ms} \\
\bottomrule
\end{tabular}
\end{table}

\begin{itemize}[leftmargin=*]
  \item \textbf{均值代替}在小集合上完全失效：用全部字段训练的网络，无法只靠三个字段做预测。
  \item \textbf{LazyVI} 精度尚可，但需要双精度，每集合 1.5—9.1 秒，全表约 25 GPU 小时，所以只跑了 200 个集合的样本。
  \item \textbf{掩码代理模型直接输出}最快，但 $M$ 误差约为本文的 4 倍。这再次说明，逐集合闭式读出是让掩码主干变准的关键。
  \item \textbf{TabPFN} 是唯一与本文精度相当的方法。
\end{itemize}

\subsection{本文与 TabPFN 分数据组看}

\fig{figures/图4.png}{127 号：本文与 TabPFN 分数据组对比\protect\newline{\footnotesize 原图：\texttt{DNN\_Aggresvation127/figures/fig2\_ours\_vs\_tabpfn.png}}}

\begin{table}[H]\centering
\caption{127 号：分数据组（本文 / TabPFN）}
\small
\begin{tabular}{lcccc}
\toprule
数据 & $V$ 偏差 & $M^{(2)}$ 误差 & $\Gamma^{(2)}$ 误差 & 关键召回 $\tau=0.7$ \\
\midrule
RTS-GMLC & $-0.003$ / $+0.004$ & 0.0085 / 0.0074 & 0.0079 / 0.0107 & 1 / 1 \\
NEM & 0.000 / $-0.018$ & \textbf{0.0144} / 0.0195 & \textbf{0.0153} / 0.0320 & 1 / 1 \\
PJM 负荷 & $-0.038$ / $+0.015$ & 0.0200 / \textbf{0.0107} & \textbf{0.0137} / 0.0162 & 0.79 / \textbf{1} \\
PJM 发电 / 联络线 & $-0.032$ / $+0.016$ & 0.0226 / \textbf{0.0161} & 0.0200 / \textbf{0.0167} & 0.25 / \textbf{1} \\
CAISO 负荷 & $-0.016$ / $+0.008$ & 0.0153 / \textbf{0.0111} & \textbf{0.0141} / 0.0146 & 1 / 1 \\
\bottomrule
\end{tabular}
\end{table}

\begin{itemize}[leftmargin=*]
  \item TabPFN 的优势集中在 \textbf{PJM}（训练行最少，本文在此有 0.03—0.04 的系统低估，PJM 发电 / 联络线在 $\tau=0.7$ 下的召回只有 0.25）；在 \textbf{NEM} 上本文明显更好。
  \item TabPFN 的平均偏差为正，说明它\textbf{有时比重训的攻击器族更强}。论文据此在局限中说明：报告的风险对拥有此类模型的攻击者只是下界，并建议把它作为攻击器族的一员。
  \item \textbf{认证之后两者没有差别}：认证前 1 为 0.959 对 0.961，前 3 为 0.979 对 0.978。
\end{itemize}

\begin{notebox}
为什么速度差距是决定性的：两者\textbf{摊销的对象不同}。TabPFN 在不同数据集之间摊销学习，每个字段集合对它都是一个新任务，要逐个目标送入 8 $\times$ 12 层 transformer；本文在同一数据集的所有字段集合之间摊销，每次查询只需几次小网络前向和三次线性求解（所有目标共用）。审计需要成千上万次查询：NEM 的 $K=2$ 全表，TabPFN 要 2.7 小时，本文 5 分钟；CAISO 规模 4 的 52,360 个集合，TabPFN 要 12 小时，本文 54 分钟。
\end{notebox}

\section{论文当前版本}

\begin{table}[H]\centering
\caption{论文版本}
\small
\begin{tabular}{lp{0.62\linewidth}}
\toprule
版本 & 内容 \\
\midrule
V5.1\_en / V5.2\_en & 估计器换成 C3、下游全部重算；估计器一节拆成“设计选择”与“外部对比”，加入 127 号的 baseline 表和图 \\
V6\_ieee & IEEE 双栏版（21 页，含 4 个附录），在 V5.2 基础上做全文文字修订 \\
IJCIP / JISA 定向稿 & 以 V6 为母版的两个投稿备选版本：IJCIP 以关键基础设施保护、物理机制为主线；JISA 以背景预算泄露画像及其验证为主线。尚未提交 \\
\bottomrule
\end{tabular}
\end{table}

\section{小结}

\begin{keybox}[小结]
\begin{itemize}[leftmargin=*]
  \item 本文估计器每集合 48—64 ms，比串行重训快 241—906 倍；快速读出与原实现逐位相同。
  \item 外部方法中只有 TabPFN 精度相当，但慢 13—110 倍；认证后两者相同。其余方法的 $M$ 误差为本文的 3—16 倍。
  \item 需要写进论文的局限：训练行少的数据集上本文会系统低估；更强的攻击器（如 TabPFN）存在时，报告的风险是下界。
\end{itemize}
\end{keybox}

\end{document}
